\begin{longtable}{|p{\textwidth}|}
\caption{Example CDL description of ' +name+ r' dataset}
rlabel{tab:cdl-'+name+r'} \\'
\hline \endhead
\hline \endfoot
netcdf latlon example\\
dimensions\\
\hline
\rowcolor{YellowGreen}  time = 1\\
\rowcolor{YellowGreen}  latitude = 360\\
\rowcolor{YellowGreen}  longitude = 720\\
\rowcolor{YellowGreen}  nv = 2\\
\hline

coordinates\\
\hline

\hangindent=1.0cm
\hangafter1
\rowcolor{Apricot}\hspace{0.5cm}int32 time (time)\\

\hangindent=1.5cm
\hangafter1
\rowcolor{Apricot}\hspace{0.5cm}\hspace{0.5cm}time:axis = "T"\\

\hangindent=1.5cm
\hangafter1
\rowcolor{Apricot}\hspace{0.5cm}\hspace{0.5cm}time:bounds = "time\_bnds"\\

\hangindent=1.5cm
\hangafter1
\rowcolor{Apricot}\hspace{0.5cm}\hspace{0.5cm}time:coverage\_content\_type = "coordinate"\\

\hangindent=1.5cm
\hangafter1
\rowcolor{Apricot}\hspace{0.5cm}\hspace{0.5cm}time:long\_name = "center time of averaging period"\\

\hangindent=1.5cm
\hangafter1
\rowcolor{Apricot}\hspace{0.5cm}\hspace{0.5cm}time:standard\_name = "time"\\

\hangindent=1.5cm
\hangafter1
\rowcolor{Apricot}\hspace{0.5cm}\hspace{0.5cm}time:units = "hours since 1992-01-01T12:00:00"\\

\hangindent=1.5cm
\hangafter1
\rowcolor{Apricot}\hspace{0.5cm}\hspace{0.5cm}time:calendar = "proleptic\_gregorian"\\

\hangindent=1.0cm
\hangafter1
\rowcolor{Apricot}\hspace{0.5cm}float32 latitude (latitude)\\

\hangindent=1.5cm
\hangafter1
\rowcolor{Apricot}\hspace{0.5cm}\hspace{0.5cm}latitude:axis = "Y"\\

\hangindent=1.5cm
\hangafter1
\rowcolor{Apricot}\hspace{0.5cm}\hspace{0.5cm}latitude:bounds = "latitude\_bnds"\\

\hangindent=1.5cm
\hangafter1
\rowcolor{Apricot}\hspace{0.5cm}\hspace{0.5cm}latitude:comment = "uniform grid spacing from -89.75 to 89.75 by 0.5"\\

\hangindent=1.5cm
\hangafter1
\rowcolor{Apricot}\hspace{0.5cm}\hspace{0.5cm}latitude:coverage\_content\_type = "coordinate"\\

\hangindent=1.5cm
\hangafter1
\rowcolor{Apricot}\hspace{0.5cm}\hspace{0.5cm}latitude:long\_name = "latitude at grid cell center"\\

\hangindent=1.5cm
\hangafter1
\rowcolor{Apricot}\hspace{0.5cm}\hspace{0.5cm}latitude:standard\_name = "latitude"\\

\hangindent=1.5cm
\hangafter1
\rowcolor{Apricot}\hspace{0.5cm}\hspace{0.5cm}latitude:units = "degrees\_north"\\

\hangindent=1.0cm
\hangafter1
\rowcolor{Apricot}\hspace{0.5cm}float32 longitude (longitude)\\

\hangindent=1.5cm
\hangafter1
\rowcolor{Apricot}\hspace{0.5cm}\hspace{0.5cm}longitude:axis = "X"\\

\hangindent=1.5cm
\hangafter1
\rowcolor{Apricot}\hspace{0.5cm}\hspace{0.5cm}longitude:bounds = "longitude\_bnds"\\

\hangindent=1.5cm
\hangafter1
\rowcolor{Apricot}\hspace{0.5cm}\hspace{0.5cm}longitude:comment = "uniform grid spacing from -179.75 to 179.75 by 0.5"\\

\hangindent=1.5cm
\hangafter1
\rowcolor{Apricot}\hspace{0.5cm}\hspace{0.5cm}longitude:coverage\_content\_type = "coordinate"\\

\hangindent=1.5cm
\hangafter1
\rowcolor{Apricot}\hspace{0.5cm}\hspace{0.5cm}longitude:long\_name = "longitude at grid cell center"\\

\hangindent=1.5cm
\hangafter1
\rowcolor{Apricot}\hspace{0.5cm}\hspace{0.5cm}longitude:standard\_name = "longitude"\\

\hangindent=1.5cm
\hangafter1
\rowcolor{Apricot}\hspace{0.5cm}\hspace{0.5cm}longitude:units = "degrees\_east"\\

\hangindent=1.0cm
\hangafter1
\rowcolor{Apricot}\hspace{0.5cm}int32 time\_bnds (time, nv)\\

\hangindent=1.5cm
\hangafter1
\rowcolor{Apricot}\hspace{0.5cm}\hspace{0.5cm}time\_bnds:comment = "Start and end times of averaging period."\\

\hangindent=1.5cm
\hangafter1
\rowcolor{Apricot}\hspace{0.5cm}\hspace{0.5cm}time\_bnds:coverage\_content\_type = "coordinate"\\

\hangindent=1.5cm
\hangafter1
\rowcolor{Apricot}\hspace{0.5cm}\hspace{0.5cm}time\_bnds:long\_name = "time bounds of averaging period"\\

\hangindent=1.0cm
\hangafter1
\rowcolor{Apricot}\hspace{0.5cm}float32 latitude\_bnds (latitude, nv)\\

\hangindent=1.5cm
\hangafter1
\rowcolor{Apricot}\hspace{0.5cm}\hspace{0.5cm}latitude\_bnds:coverage\_content\_type = "coordinate"\\

\hangindent=1.5cm
\hangafter1
\rowcolor{Apricot}\hspace{0.5cm}\hspace{0.5cm}latitude\_bnds:long\_name = "latitude bounds grid cells"\\

\hangindent=1.0cm
\hangafter1
\rowcolor{Apricot}\hspace{0.5cm}float32 longitude\_bnds (longitude, nv)\\

\hangindent=1.5cm
\hangafter1
\rowcolor{Apricot}\hspace{0.5cm}\hspace{0.5cm}longitude\_bnds:coverage\_content\_type = "coordinate"\\

\hangindent=1.5cm
\hangafter1
\rowcolor{Apricot}\hspace{0.5cm}\hspace{0.5cm}longitude\_bnds:long\_name = "longitude bounds grid cells"\\
\hline

data variables\\
\hline
\\
\hangindent=1.0cm
\hangafter1
\hspace{0.5cm}float32 OBP (time, latitude, longitude)\\
\\
\hangindent=1.5cm
\hangafter1
\hspace{0.5cm}\hspace{0.5cm}OBP:\_FillValue = "9.969209968386869e+36"\\
\\
\hangindent=1.5cm
\hangafter1
\hspace{0.5cm}\hspace{0.5cm}OBP:coverage\_content\_type = "modelResult"\\
\\
\hangindent=1.5cm
\hangafter1
\hspace{0.5cm}\hspace{0.5cm}OBP:long\_name = "Ocean bottom pressure given as equivalent water thickness"\\
\\
\hangindent=1.5cm
\hangafter1
\hspace{0.5cm}\hspace{0.5cm}OBP:units = "m"\\
\\
\hangindent=1.5cm
\hangafter1
\hspace{0.5cm}\hspace{0.5cm}OBP:comment = "OBP excludes the contribution from global mean atmospheric pressure and is therefore suitable for comparisons with GRACE data products. OBP is calculated as follows. First, we calculate ocean hydrostatic bottom pressure anomaly, PHIBOT, with PHIBOT = p\_b/rhoConst - gH(t), where p\_b = model ocean hydrostatic bottom pressure, rhoConst = reference density (1029 kg m-3), g is acceleration due to gravity (9.81 m s-2), and H(t) is model depth at time t. Then, OBP = PHIBOT/g + corrections for i) global mean steric sea level changes related to density changes in the Boussinesq volume-conserving model (Greatbatch correction, see sterGloH) and ii) global mean atmospheric pressure variations. Use OBP for comparisons with ocean bottom pressure data products that have been corrected for global mean atmospheric pressure variations. GRACE data typically ARE corrected for global mean atmospheric pressure variations. In contrast, ocean bottom pressure gauge data typically ARE NOT corrected for global mean atmospheric pressure variations."\\
\\
\hangindent=1.5cm
\hangafter1
\hspace{0.5cm}\hspace{0.5cm}OBP:coordinates = "time"\\
\\
\hangindent=1.5cm
\hangafter1
\hspace{0.5cm}\hspace{0.5cm}OBP:valid\_min = "-2.544442892074585"\\
\\
\hangindent=1.5cm
\hangafter1
\hspace{0.5cm}\hspace{0.5cm}OBP:valid\_max = "72.1243667602539"\\
\\
\hangindent=1.0cm
\hangafter1
\hspace{0.5cm}float32 OBPGMAP (time, latitude, longitude)\\
\\
\hangindent=1.5cm
\hangafter1
\hspace{0.5cm}\hspace{0.5cm}OBPGMAP:\_FillValue = "9.969209968386869e+36"\\
\\
\hangindent=1.5cm
\hangafter1
\hspace{0.5cm}\hspace{0.5cm}OBPGMAP:coverage\_content\_type = "modelResult"\\
\\
\hangindent=1.5cm
\hangafter1
\hspace{0.5cm}\hspace{0.5cm}OBPGMAP:long\_name = "Ocean bottom pressure given as equivalent water thickness, includes global mean atmospheric pressure"\\
\\
\hangindent=1.5cm
\hangafter1
\hspace{0.5cm}\hspace{0.5cm}OBPGMAP:units = "m"\\
\\
\hangindent=1.5cm
\hangafter1
\hspace{0.5cm}\hspace{0.5cm}OBPGMAP:comment = "OBPGMAP includes the contribution from global mean atmospheric pressure and is therefore suitable for comparisons with ocean bottom pressure gauge data products. OBPGMAP is calculated as follows. First, we calculate ocean hydrostatic bottom pressure anomaly, PHIBOT, with PHIBOT = p\_b/rhoConst - gH(t), where p\_b = model ocean hydrostatic bottom pressure, rhoConst = reference density (1029 kg m-3), g is acceleration due to gravity (9.81 m s-2), and H(t) is model depth at time t. Then, OBPGMAP= PHIBOT/g + corrections for global mean steric sea level changes related to density changes in the Boussinesq volume-conserving model (Greatbatch correction, see sterGloH). Use OBPGMAP for comparisons with ocean bottom pressure data products that have NOT been corrected for global mean atmospheric pressure variations. GRACE data typically ARE corrected for global mean atmospheric pressure variations. In contrast, ocean bottom pressure gauge data typically ARE NOT corrected for global mean atmospheric pressure variations."\\
\\
\hangindent=1.5cm
\hangafter1
\hspace{0.5cm}\hspace{0.5cm}OBPGMAP:coordinates = "time"\\
\\
\hangindent=1.5cm
\hangafter1
\hspace{0.5cm}\hspace{0.5cm}OBPGMAP:valid\_min = "7.395928859710693"\\
\\
\hangindent=1.5cm
\hangafter1
\hspace{0.5cm}\hspace{0.5cm}OBPGMAP:valid\_max = "82.14805603027344"\\
\hline
\end{longtable}
netcdf native example\\
dimensions\\
\hline
\rowcolor{YellowGreen}  tile = 13\\
\rowcolor{YellowGreen}  j = 90\\
\rowcolor{YellowGreen}  i = 90\\
\rowcolor{YellowGreen}  j\_g = 90\\
\rowcolor{YellowGreen}  i\_g = 90\\
\rowcolor{YellowGreen}  time = 1\\
\rowcolor{YellowGreen}  nv = 2\\
\rowcolor{YellowGreen}  nb = 4\\
\hline

coordinates\\
\hline

\hangindent=1.0cm
\hangafter1
\rowcolor{Apricot}\hspace{0.5cm}int32 i (i)\\

\hangindent=1.5cm
\hangafter1
\rowcolor{Apricot}\hspace{0.5cm}\hspace{0.5cm}i:axis = "X"\\

\hangindent=1.5cm
\hangafter1
\rowcolor{Apricot}\hspace{0.5cm}\hspace{0.5cm}i:long\_name = "grid index in x for variables at tracer and 'v' locations"\\

\hangindent=1.5cm
\hangafter1
\rowcolor{Apricot}\hspace{0.5cm}\hspace{0.5cm}i:swap\_dim = "XC"\\

\hangindent=1.5cm
\hangafter1
\rowcolor{Apricot}\hspace{0.5cm}\hspace{0.5cm}i:comment = "In the Arakawa C-grid system, tracer (e.g., THETA) and 'v' variables (e.g., VVEL) have the same x coordinate on the model grid."\\

\hangindent=1.5cm
\hangafter1
\rowcolor{Apricot}\hspace{0.5cm}\hspace{0.5cm}i:coverage\_content\_type = "coordinate"\\

\hangindent=1.0cm
\hangafter1
\rowcolor{Apricot}\hspace{0.5cm}int32 i\_g (i\_g)\\

\hangindent=1.5cm
\hangafter1
\rowcolor{Apricot}\hspace{0.5cm}\hspace{0.5cm}i\_g:axis = "X"\\

\hangindent=1.5cm
\hangafter1
\rowcolor{Apricot}\hspace{0.5cm}\hspace{0.5cm}i\_g:long\_name = "grid index in x for variables at 'u' and 'g' locations"\\

\hangindent=1.5cm
\hangafter1
\rowcolor{Apricot}\hspace{0.5cm}\hspace{0.5cm}i\_g:c\_grid\_axis\_shift = "-0.5"\\

\hangindent=1.5cm
\hangafter1
\rowcolor{Apricot}\hspace{0.5cm}\hspace{0.5cm}i\_g:swap\_dim = "XG"\\

\hangindent=1.5cm
\hangafter1
\rowcolor{Apricot}\hspace{0.5cm}\hspace{0.5cm}i\_g:comment = "In the Arakawa C-grid system, 'u' (e.g., UVEL) and 'g' variables (e.g., XG) have the same x coordinate on the model grid."\\

\hangindent=1.5cm
\hangafter1
\rowcolor{Apricot}\hspace{0.5cm}\hspace{0.5cm}i\_g:coverage\_content\_type = "coordinate"\\

\hangindent=1.0cm
\hangafter1
\rowcolor{Apricot}\hspace{0.5cm}int32 j (j)\\

\hangindent=1.5cm
\hangafter1
\rowcolor{Apricot}\hspace{0.5cm}\hspace{0.5cm}j:axis = "Y"\\

\hangindent=1.5cm
\hangafter1
\rowcolor{Apricot}\hspace{0.5cm}\hspace{0.5cm}j:long\_name = "grid index in y for variables at tracer and 'u' locations"\\

\hangindent=1.5cm
\hangafter1
\rowcolor{Apricot}\hspace{0.5cm}\hspace{0.5cm}j:swap\_dim = "YC"\\

\hangindent=1.5cm
\hangafter1
\rowcolor{Apricot}\hspace{0.5cm}\hspace{0.5cm}j:comment = "In the Arakawa C-grid system, tracer (e.g., THETA) and 'u' variables (e.g., UVEL) have the same y coordinate on the model grid."\\

\hangindent=1.5cm
\hangafter1
\rowcolor{Apricot}\hspace{0.5cm}\hspace{0.5cm}j:coverage\_content\_type = "coordinate"\\

\hangindent=1.0cm
\hangafter1
\rowcolor{Apricot}\hspace{0.5cm}int32 j\_g (j\_g)\\

\hangindent=1.5cm
\hangafter1
\rowcolor{Apricot}\hspace{0.5cm}\hspace{0.5cm}j\_g:axis = "Y"\\

\hangindent=1.5cm
\hangafter1
\rowcolor{Apricot}\hspace{0.5cm}\hspace{0.5cm}j\_g:long\_name = "grid index in y for variables at 'v' and 'g' locations"\\

\hangindent=1.5cm
\hangafter1
\rowcolor{Apricot}\hspace{0.5cm}\hspace{0.5cm}j\_g:c\_grid\_axis\_shift = "-0.5"\\

\hangindent=1.5cm
\hangafter1
\rowcolor{Apricot}\hspace{0.5cm}\hspace{0.5cm}j\_g:swap\_dim = "YG"\\

\hangindent=1.5cm
\hangafter1
\rowcolor{Apricot}\hspace{0.5cm}\hspace{0.5cm}j\_g:comment = "In the Arakawa C-grid system, 'v' (e.g., VVEL) and 'g' variables (e.g., XG) have the same y coordinate."\\

\hangindent=1.5cm
\hangafter1
\rowcolor{Apricot}\hspace{0.5cm}\hspace{0.5cm}j\_g:coverage\_content\_type = "coordinate"\\

\hangindent=1.0cm
\hangafter1
\rowcolor{Apricot}\hspace{0.5cm}int32 tile (tile)\\

\hangindent=1.5cm
\hangafter1
\rowcolor{Apricot}\hspace{0.5cm}\hspace{0.5cm}tile:long\_name = "lat-lon-cap tile index"\\

\hangindent=1.5cm
\hangafter1
\rowcolor{Apricot}\hspace{0.5cm}\hspace{0.5cm}tile:comment = "The ECCO V4 horizontal model grid is divided into 13 tiles of 90x90 cells for convenience."\\

\hangindent=1.5cm
\hangafter1
\rowcolor{Apricot}\hspace{0.5cm}\hspace{0.5cm}tile:coverage\_content\_type = "coordinate"\\

\hangindent=1.0cm
\hangafter1
\rowcolor{Apricot}\hspace{0.5cm}int32 time (time)\\

\hangindent=1.5cm
\hangafter1
\rowcolor{Apricot}\hspace{0.5cm}\hspace{0.5cm}time:long\_name = "center time of averaging period"\\

\hangindent=1.5cm
\hangafter1
\rowcolor{Apricot}\hspace{0.5cm}\hspace{0.5cm}time:axis = "T"\\

\hangindent=1.5cm
\hangafter1
\rowcolor{Apricot}\hspace{0.5cm}\hspace{0.5cm}time:bounds = "time\_bnds"\\

\hangindent=1.5cm
\hangafter1
\rowcolor{Apricot}\hspace{0.5cm}\hspace{0.5cm}time:coverage\_content\_type = "coordinate"\\

\hangindent=1.5cm
\hangafter1
\rowcolor{Apricot}\hspace{0.5cm}\hspace{0.5cm}time:standard\_name = "time"\\

\hangindent=1.5cm
\hangafter1
\rowcolor{Apricot}\hspace{0.5cm}\hspace{0.5cm}time:units = "hours since 1992-01-01T12:00:00"\\

\hangindent=1.5cm
\hangafter1
\rowcolor{Apricot}\hspace{0.5cm}\hspace{0.5cm}time:calendar = "proleptic\_gregorian"\\

\hangindent=1.0cm
\hangafter1
\rowcolor{Apricot}\hspace{0.5cm}float32 XC (tile, j, i)\\

\hangindent=1.5cm
\hangafter1
\rowcolor{Apricot}\hspace{0.5cm}\hspace{0.5cm}XC:long\_name = "longitude of tracer grid cell center"\\

\hangindent=1.5cm
\hangafter1
\rowcolor{Apricot}\hspace{0.5cm}\hspace{0.5cm}XC:units = "degrees\_east"\\

\hangindent=1.5cm
\hangafter1
\rowcolor{Apricot}\hspace{0.5cm}\hspace{0.5cm}XC:coordinate = "YC XC"\\

\hangindent=1.5cm
\hangafter1
\rowcolor{Apricot}\hspace{0.5cm}\hspace{0.5cm}XC:bounds = "XC\_bnds"\\

\hangindent=1.5cm
\hangafter1
\rowcolor{Apricot}\hspace{0.5cm}\hspace{0.5cm}XC:comment = "nonuniform grid spacing"\\

\hangindent=1.5cm
\hangafter1
\rowcolor{Apricot}\hspace{0.5cm}\hspace{0.5cm}XC:coverage\_content\_type = "coordinate"\\

\hangindent=1.5cm
\hangafter1
\rowcolor{Apricot}\hspace{0.5cm}\hspace{0.5cm}XC:standard\_name = "longitude"\\

\hangindent=1.0cm
\hangafter1
\rowcolor{Apricot}\hspace{0.5cm}float32 YC (tile, j, i)\\

\hangindent=1.5cm
\hangafter1
\rowcolor{Apricot}\hspace{0.5cm}\hspace{0.5cm}YC:long\_name = "latitude of tracer grid cell center"\\

\hangindent=1.5cm
\hangafter1
\rowcolor{Apricot}\hspace{0.5cm}\hspace{0.5cm}YC:units = "degrees\_north"\\

\hangindent=1.5cm
\hangafter1
\rowcolor{Apricot}\hspace{0.5cm}\hspace{0.5cm}YC:coordinate = "YC XC"\\

\hangindent=1.5cm
\hangafter1
\rowcolor{Apricot}\hspace{0.5cm}\hspace{0.5cm}YC:bounds = "YC\_bnds"\\

\hangindent=1.5cm
\hangafter1
\rowcolor{Apricot}\hspace{0.5cm}\hspace{0.5cm}YC:comment = "nonuniform grid spacing"\\

\hangindent=1.5cm
\hangafter1
\rowcolor{Apricot}\hspace{0.5cm}\hspace{0.5cm}YC:coverage\_content\_type = "coordinate"\\

\hangindent=1.5cm
\hangafter1
\rowcolor{Apricot}\hspace{0.5cm}\hspace{0.5cm}YC:standard\_name = "latitude"\\

\hangindent=1.0cm
\hangafter1
\rowcolor{Apricot}\hspace{0.5cm}float32 XG (tile, j\_g, i\_g)\\

\hangindent=1.5cm
\hangafter1
\rowcolor{Apricot}\hspace{0.5cm}\hspace{0.5cm}XG:long\_name = "longitude of 'southwest' corner of tracer grid cell"\\

\hangindent=1.5cm
\hangafter1
\rowcolor{Apricot}\hspace{0.5cm}\hspace{0.5cm}XG:units = "degrees\_east"\\

\hangindent=1.5cm
\hangafter1
\rowcolor{Apricot}\hspace{0.5cm}\hspace{0.5cm}XG:coordinate = "YG XG"\\

\hangindent=1.5cm
\hangafter1
\rowcolor{Apricot}\hspace{0.5cm}\hspace{0.5cm}XG:comment = "Nonuniform grid spacing. Note: 'southwest' does not correspond to geographic orientation but is used for convenience to describe the computational grid. See MITgcm dcoumentation for details."\\

\hangindent=1.5cm
\hangafter1
\rowcolor{Apricot}\hspace{0.5cm}\hspace{0.5cm}XG:coverage\_content\_type = "coordinate"\\

\hangindent=1.5cm
\hangafter1
\rowcolor{Apricot}\hspace{0.5cm}\hspace{0.5cm}XG:standard\_name = "longitude"\\

\hangindent=1.0cm
\hangafter1
\rowcolor{Apricot}\hspace{0.5cm}float32 YG (tile, j\_g, i\_g)\\

\hangindent=1.5cm
\hangafter1
\rowcolor{Apricot}\hspace{0.5cm}\hspace{0.5cm}YG:long\_name = "latitude of 'southwest' corner of tracer grid cell"\\

\hangindent=1.5cm
\hangafter1
\rowcolor{Apricot}\hspace{0.5cm}\hspace{0.5cm}YG:units = "degrees\_north"\\

\hangindent=1.5cm
\hangafter1
\rowcolor{Apricot}\hspace{0.5cm}\hspace{0.5cm}YG:coordinate = "YG XG"\\

\hangindent=1.5cm
\hangafter1
\rowcolor{Apricot}\hspace{0.5cm}\hspace{0.5cm}YG:comment = "Nonuniform grid spacing. Note: 'southwest' does not correspond to geographic orientation but is used for convenience to describe the computational grid. See MITgcm dcoumentation for details."\\

\hangindent=1.5cm
\hangafter1
\rowcolor{Apricot}\hspace{0.5cm}\hspace{0.5cm}YG:coverage\_content\_type = "coordinate"\\

\hangindent=1.5cm
\hangafter1
\rowcolor{Apricot}\hspace{0.5cm}\hspace{0.5cm}YG:standard\_name = "latitude"\\

\hangindent=1.0cm
\hangafter1
\rowcolor{Apricot}\hspace{0.5cm}int32 time\_bnds (time, nv)\\

\hangindent=1.5cm
\hangafter1
\rowcolor{Apricot}\hspace{0.5cm}\hspace{0.5cm}time\_bnds:comment = "Start and end times of averaging period."\\

\hangindent=1.5cm
\hangafter1
\rowcolor{Apricot}\hspace{0.5cm}\hspace{0.5cm}time\_bnds:coverage\_content\_type = "coordinate"\\

\hangindent=1.5cm
\hangafter1
\rowcolor{Apricot}\hspace{0.5cm}\hspace{0.5cm}time\_bnds:long\_name = "time bounds of averaging period"\\

\hangindent=1.0cm
\hangafter1
\rowcolor{Apricot}\hspace{0.5cm}float32 XC\_bnds (tile, j, i, nb)\\

\hangindent=1.5cm
\hangafter1
\rowcolor{Apricot}\hspace{0.5cm}\hspace{0.5cm}XC\_bnds:comment = "Bounds array follows CF conventions. XC\_bnds[i,j,0] = 'southwest' corner (j-1, i-1), XC\_bnds[i,j,1] = 'southeast' corner (j-1, i+1), XC\_bnds[i,j,2] = 'northeast' corner (j+1, i+1), XC\_bnds[i,j,3]  = 'northwest' corner (j+1, i-1). Note: 'southwest', 'southeast', northwest', and 'northeast' do not correspond to geographic orientation but are used for convenience to describe the computational grid. See MITgcm dcoumentation for details."\\

\hangindent=1.5cm
\hangafter1
\rowcolor{Apricot}\hspace{0.5cm}\hspace{0.5cm}XC\_bnds:coverage\_content\_type = "coordinate"\\

\hangindent=1.5cm
\hangafter1
\rowcolor{Apricot}\hspace{0.5cm}\hspace{0.5cm}XC\_bnds:long\_name = "longitudes of tracer grid cell corners"\\

\hangindent=1.0cm
\hangafter1
\rowcolor{Apricot}\hspace{0.5cm}float32 YC\_bnds (tile, j, i, nb)\\

\hangindent=1.5cm
\hangafter1
\rowcolor{Apricot}\hspace{0.5cm}\hspace{0.5cm}YC\_bnds:comment = "Bounds array follows CF conventions. YC\_bnds[i,j,0] = 'southwest' corner (j-1, i-1), YC\_bnds[i,j,1] = 'southeast' corner (j-1, i+1), YC\_bnds[i,j,2] = 'northeast' corner (j+1, i+1), YC\_bnds[i,j,3]  = 'northwest' corner (j+1, i-1). Note: 'southwest', 'southeast', northwest', and 'northeast' do not correspond to geographic orientation but are used for convenience to describe the computational grid. See MITgcm dcoumentation for details."\\

\hangindent=1.5cm
\hangafter1
\rowcolor{Apricot}\hspace{0.5cm}\hspace{0.5cm}YC\_bnds:coverage\_content\_type = "coordinate"\\

\hangindent=1.5cm
\hangafter1
\rowcolor{Apricot}\hspace{0.5cm}\hspace{0.5cm}YC\_bnds:long\_name = "latitudes of tracer grid cell corners"\\
\hline

data variables\\
\hline
\\
\hangindent=1.0cm
\hangafter1
\hspace{0.5cm}float32 MXLDEPTH (time, tile, j, i)\\
\\
\hangindent=1.5cm
\hangafter1
\hspace{0.5cm}\hspace{0.5cm}MXLDEPTH:\_FillValue = "9.969209968386869e+36"\\
\\
\hangindent=1.5cm
\hangafter1
\hspace{0.5cm}\hspace{0.5cm}MXLDEPTH:long\_name = "Mixed-layer depth diagnosed using the temperature difference criterion of Kara et al., 2000"\\
\\
\hangindent=1.5cm
\hangafter1
\hspace{0.5cm}\hspace{0.5cm}MXLDEPTH:units = "m"\\
\\
\hangindent=1.5cm
\hangafter1
\hspace{0.5cm}\hspace{0.5cm}MXLDEPTH:coverage\_content\_type = "modelResult"\\
\\
\hangindent=1.5cm
\hangafter1
\hspace{0.5cm}\hspace{0.5cm}MXLDEPTH:standard\_name = "ocean\_mixed\_layer\_thickness"\\
\\
\hangindent=1.5cm
\hangafter1
\hspace{0.5cm}\hspace{0.5cm}MXLDEPTH:comment = "Mixed-layer depth as determined by the depth where waters are first 0.8 degrees Celsius colder than the surface. See Kara et al. (JGR, 2000). . Note: the Kara et al. criterion may not be appropriate for some applications. If needed, mixed layer depth can be calculated using different criteria. See vertical density stratification (DRHODR) and density anomaly (RHOAnoma)."\\
\\
\hangindent=1.5cm
\hangafter1
\hspace{0.5cm}\hspace{0.5cm}MXLDEPTH:coordinates = "time XC YC"\\
\\
\hangindent=1.5cm
\hangafter1
\hspace{0.5cm}\hspace{0.5cm}MXLDEPTH:valid\_min = "5.000001430511475"\\
\\
\hangindent=1.5cm
\hangafter1
\hspace{0.5cm}\hspace{0.5cm}MXLDEPTH:valid\_max = "5331.2001953125"\\
\hline
\end{longtable}